\section{洪熙至天顺的补充编年节点（二）}
以下节点继续沿同一编年脉络展开。每一锚点均紧接可供核查的事件叙述；来源能确认何种行动、文书或变动，不能自动消除不同记录之间的立场差异。

\subsection{1458年前后的续接}

\eventanchor{ML-1458-003}{明·1458（天顺二年）}
\paragraph{1458（天顺二年）：地方灾异、赈济与水利条目为本年社会史提供入口，地方志的修纂层次需要说明。} 这项记录把本章所见的制度运作落实到具体年份。其形成背景是：自然风险与地方报告促成朝廷或地方的应对记录。 它带来的、而且仍可由后续材料追踪的改变是：赈济、迁徙和损失的实际范围仍须以受灾地区材料分别核验。 《英宗实录》天顺二年条；《明史·五行志》；受灾府县地方志与奏疏 提供相互检验的起点；但桥接条目只标示可回查的年度问题与材料链；实录有中央选择性，崇祯期尤其需要奏疏、地方、朝鲜及满文材料交叉。 因此本段仅据材料确认事件的时序、行动和可见后果，并把政策文本与地方实践、报称数字和社会经历分开处理。

\eventanchor{MM-1458-CROWN-PRINCE-RESTORED}{明·1458（天顺二年）}
\paragraph{1458（天顺二年）：英宗复立朱见深为皇太子} 这项记录把本章所见的制度运作落实到具体年份。其形成背景是：朱见济已死，英宗复位后须恢复本支继承次序。 它带来的、而且仍可由后续材料追踪的改变是：朱见深重新成为法定继承人，为1464年宪宗即位预置了宫廷与礼制程序。 《英宗实录》天顺二年条；《明史·宪宗本纪》；《明史·英宗后妃传》 提供相互检验的起点；但复立事与诏令可靠；朝臣对景泰旧臣处置和储位安排的影响须结合奏疏辨析。 因此本段仅据材料确认事件的时序、行动和可见后果，并把政策文本与地方实践、报称数字和社会经历分开处理。

\eventanchor{ML-1459-001}{明·1459（天顺三年）}
\paragraph{1459（天顺三年）：景泰、天顺时期的继承、人事与诏令链条可在本年材料中定位，政变责任不可由单一史书裁断。} 这项记录把本章所见的制度运作落实到具体年份。其形成背景是：继承、官僚协调或皇权运作的压力使问题进入朝廷程序。 它带来的、而且仍可由后续材料追踪的改变是：它影响后续人事与政策议程；动机和派系标签不能由单一叙事裁定。 《英宗实录》天顺三年条；《明史·本纪》及相关列传；诏令、奏疏或会典版本 提供相互检验的起点；但桥接条目只标示可回查的年度问题与材料链；实录有中央选择性，崇祯期尤其需要奏疏、地方、朝鲜及满文材料交叉。 因此本段仅据材料确认事件的时序、行动和可见后果，并把政策文本与地方实践、报称数字和社会经历分开处理。

\eventanchor{ML-1459-002}{明·1459（天顺三年）}
\paragraph{1459（天顺三年）：蒙古诸部、朝鲜及西北绿洲的交涉在本年留下多方材料，应以具体使团和地名校正概称。} 这项记录把本章所见的制度运作落实到具体年份。其形成背景是：朝贡名义、边境安全与港口贸易的利益使交涉持续发生。 它带来的、而且仍可由后续材料追踪的改变是：它为后续册封、互市或海防安排留下文书与跨政权比较线索。 《英宗实录》天顺三年条；《明史·外国传》；《朝鲜王朝实录》、琉球《历代宝案》或海外档案 提供相互检验的起点；但桥接条目只标示可回查的年度问题与材料链；实录有中央选择性，崇祯期尤其需要奏疏、地方、朝鲜及满文材料交叉。 因此本段仅据材料确认事件的时序、行动和可见后果，并把政策文本与地方实践、报称数字和社会经历分开处理。

\eventanchor{ML-1459-003}{明·1459（天顺三年）}
\paragraph{1459（天顺三年）：科举、学校和法律条文在本年制度材料中可追踪，不能据规范文本推定州县实践一致。} 这项记录把本章所见的制度运作落实到具体年份。其形成背景是：制度、教育或知识传播的需要推动了相关行动或文本形成。 它带来的、而且仍可由后续材料追踪的改变是：它提供制度和传播史的锚点，不能据此推断社会整体接受。 《英宗实录》天顺三年条；《明史·选举志》或《艺文志》；文集、碑刻或书目材料 提供相互检验的起点；但桥接条目只标示可回查的年度问题与材料链；实录有中央选择性，崇祯期尤其需要奏疏、地方、朝鲜及满文材料交叉。 因此本段仅据材料确认事件的时序、行动和可见后果，并把政策文本与地方实践、报称数字和社会经历分开处理。

\eventanchor{ML-1459-004}{明·1459（天顺三年）}
\paragraph{1459（天顺三年）：土木危机后的边镇整顿、京师防务或复辟前后军政调度在本年材料中构成连续问题。} 这项记录把本章所见的制度运作落实到具体年份。其形成背景是：前期边防、地方武装或跨区域资源调动的压力推动了该行动。 它带来的、而且仍可由后续材料追踪的改变是：它改变了后续调兵、补给或地方秩序的条件；战果和伤亡须分开核对。 《英宗实录》天顺三年条；《明史·兵志》及相关列传；边镇、地方志或对方材料 提供相互检验的起点；但桥接条目只标示可回查的年度问题与材料链；实录有中央选择性，崇祯期尤其需要奏疏、地方、朝鲜及满文材料交叉。 因此本段仅据材料确认事件的时序、行动和可见后果，并把政策文本与地方实践、报称数字和社会经历分开处理。

\eventanchor{ML-1460-001}{明·1460（天顺四年）}
\paragraph{1460（天顺四年）：科举、学校和法律条文在本年制度材料中可追踪，不能据规范文本推定州县实践一致。} 这项记录把本章所见的制度运作落实到具体年份。其形成背景是：制度、教育或知识传播的需要推动了相关行动或文本形成。 它带来的、而且仍可由后续材料追踪的改变是：它提供制度和传播史的锚点，不能据此推断社会整体接受。 《英宗实录》天顺四年条；《明史·选举志》或《艺文志》；文集、碑刻或书目材料 提供相互检验的起点；但桥接条目只标示可回查的年度问题与材料链；实录有中央选择性，崇祯期尤其需要奏疏、地方、朝鲜及满文材料交叉。 因此本段仅据材料确认事件的时序、行动和可见后果，并把政策文本与地方实践、报称数字和社会经历分开处理。

\eventanchor{ML-1460-002}{明·1460（天顺四年）}
\paragraph{1460（天顺四年）：地方灾异、赈济与水利条目为本年社会史提供入口，地方志的修纂层次需要说明。} 这项记录把本章所见的制度运作落实到具体年份。其形成背景是：自然风险与地方报告促成朝廷或地方的应对记录。 它带来的、而且仍可由后续材料追踪的改变是：赈济、迁徙和损失的实际范围仍须以受灾地区材料分别核验。 《英宗实录》天顺四年条；《明史·五行志》；受灾府县地方志与奏疏 提供相互检验的起点；但桥接条目只标示可回查的年度问题与材料链；实录有中央选择性，崇祯期尤其需要奏疏、地方、朝鲜及满文材料交叉。 因此本段仅据材料确认事件的时序、行动和可见后果，并把政策文本与地方实践、报称数字和社会经历分开处理。

\eventanchor{ML-1460-003}{明·1460（天顺四年）}
\paragraph{1460（天顺四年）：景泰、天顺时期的继承、人事与诏令链条可在本年材料中定位，政变责任不可由单一史书裁断。} 这项记录把本章所见的制度运作落实到具体年份。其形成背景是：继承、官僚协调或皇权运作的压力使问题进入朝廷程序。 它带来的、而且仍可由后续材料追踪的改变是：它影响后续人事与政策议程；动机和派系标签不能由单一叙事裁定。 《英宗实录》天顺四年条；《明史·本纪》及相关列传；诏令、奏疏或会典版本 提供相互检验的起点；但桥接条目只标示可回查的年度问题与材料链；实录有中央选择性，崇祯期尤其需要奏疏、地方、朝鲜及满文材料交叉。 因此本段仅据材料确认事件的时序、行动和可见后果，并把政策文本与地方实践、报称数字和社会经历分开处理。

\eventanchor{MM-1460-SHI-HENG-FALL}{明·1460（天顺四年）}
\paragraph{1460（天顺四年）：石亨被逮下狱，复辟功臣集团失去中枢优势} 这项记录把本章所见的制度运作落实到具体年份。其形成背景是：石亨在夺门后掌兵并广植私人，已与文官及皇帝形成权力摩擦。 它带来的、而且仍可由后续材料追踪的改变是：英宗削弱了复辟功臣的军权；曹吉祥及其养子集团的紧张关系仍未解除。 《英宗实录》天顺四年条；《明史·石亨传》；《明史·曹吉祥传》 提供相互检验的起点；但石亨被逮及其后处分可靠；“结党谋反”等案情须与复辟后政治清算的语境一并阅读。 因此本段仅据材料确认事件的时序、行动和可见后果，并把政策文本与地方实践、报称数字和社会经历分开处理。

\eventanchor{ML-1461-001}{明·1461（天顺五年）八月7日}
\paragraph{1461（天顺五年）八月7日：曹钦之乱：曹钦叛乱并试图攻入北京，但被捕并被迫自杀} 这项记录把本章所见的制度运作落实到具体年份。其形成背景是：前期边防、地方武装或跨区域资源调动的压力推动了该行动。 它带来的、而且仍可由后续材料追踪的改变是：它改变了后续调兵、补给或地方秩序的条件；战果和伤亡须分开核对。 《英宗实录》天顺五年条；《明史·兵志》及相关列传；边镇、地方志或对方材料 提供相互检验的起点；但译写候选以编年材料定位；实录记载反映朝廷选择，崇祯期须另以奏疏、地方及敌对政权材料交叉。 因此本段仅据材料确认事件的时序、行动和可见后果，并把政策文本与地方实践、报称数字和社会经历分开处理。

\eventanchor{ML-1461-002}{明·1461（天顺五年）}
\paragraph{1461（天顺五年）：土木危机后的边镇整顿、京师防务或复辟前后军政调度在本年材料中构成连续问题。} 这项记录把本章所见的制度运作落实到具体年份。其形成背景是：前期边防、地方武装或跨区域资源调动的压力推动了该行动。 它带来的、而且仍可由后续材料追踪的改变是：它改变了后续调兵、补给或地方秩序的条件；战果和伤亡须分开核对。 《英宗实录》天顺五年条；《明史·兵志》及相关列传；边镇、地方志或对方材料 提供相互检验的起点；但桥接条目只标示可回查的年度问题与材料链；实录有中央选择性，崇祯期尤其需要奏疏、地方、朝鲜及满文材料交叉。 因此本段仅据材料确认事件的时序、行动和可见后果，并把政策文本与地方实践、报称数字和社会经历分开处理。

\eventanchor{ML-1461-003}{明·1461（天顺五年）}
\paragraph{1461（天顺五年）：军饷、漕运和户籍赋役的本年记录揭示恢复秩序的行政成本，账面额与实收须分列。} 这项记录把本章所见的制度运作落实到具体年份。其形成背景是：财政征解、交通工程或货币支付的压力使相关措施进入政务。 它带来的、而且仍可由后续材料追踪的改变是：其法定安排不等于地方执行，后续须以地域材料追踪负担与成效。 《英宗实录》天顺五年条；《明会典》相关条目；地方志、黄册/契约/河工材料 提供相互检验的起点；但桥接条目只标示可回查的年度问题与材料链；实录有中央选择性，崇祯期尤其需要奏疏、地方、朝鲜及满文材料交叉。 因此本段仅据材料确认事件的时序、行动和可见后果，并把政策文本与地方实践、报称数字和社会经历分开处理。

\eventanchor{MM-1461-CAO-QIN-REBELLION}{明·1461（天顺五年）七月}
\paragraph{1461（天顺五年）七月：曹钦在北京发动兵变，攻东安门等处后被官军平定} 这项记录把本章所见的制度运作落实到具体年份。其形成背景是：曹吉祥、曹钦集团在石亨失势后担忧被清算，且仍掌握部分京营关系。 它带来的、而且仍可由后续材料追踪的改变是：朝廷加强对宦官和武臣私属的防范；宫城、城门与京营的控制成为后续中枢安全议题。 《英宗实录》天顺五年七月条；《明史·曹吉祥传》；《明史·孙镗传》 提供相互检验的起点；但兵变的发生、主要地点和结局可靠；参战人数、城内破坏及每名将领责任多来自朝廷战报。 因此本段仅据材料确认事件的时序、行动和可见后果，并把政策文本与地方实践、报称数字和社会经历分开处理。

\eventanchor{ML-1462-001}{明·1462（天顺六年）}
\paragraph{1462（天顺六年）：军饷、漕运和户籍赋役的本年记录揭示恢复秩序的行政成本，账面额与实收须分列。} 这项记录把本章所见的制度运作落实到具体年份。其形成背景是：财政征解、交通工程或货币支付的压力使相关措施进入政务。 它带来的、而且仍可由后续材料追踪的改变是：其法定安排不等于地方执行，后续须以地域材料追踪负担与成效。 《英宗实录》天顺六年条；《明会典》相关条目；地方志、黄册/契约/河工材料 提供相互检验的起点；但桥接条目只标示可回查的年度问题与材料链；实录有中央选择性，崇祯期尤其需要奏疏、地方、朝鲜及满文材料交叉。 因此本段仅据材料确认事件的时序、行动和可见后果，并把政策文本与地方实践、报称数字和社会经历分开处理。

\eventanchor{ML-1462-002}{明·1462（天顺六年）}
\paragraph{1462（天顺六年）：地方灾异、赈济与水利条目为本年社会史提供入口，地方志的修纂层次需要说明。} 这项记录把本章所见的制度运作落实到具体年份。其形成背景是：自然风险与地方报告促成朝廷或地方的应对记录。 它带来的、而且仍可由后续材料追踪的改变是：赈济、迁徙和损失的实际范围仍须以受灾地区材料分别核验。 《英宗实录》天顺六年条；《明史·五行志》；受灾府县地方志与奏疏 提供相互检验的起点；但桥接条目只标示可回查的年度问题与材料链；实录有中央选择性，崇祯期尤其需要奏疏、地方、朝鲜及满文材料交叉。 因此本段仅据材料确认事件的时序、行动和可见后果，并把政策文本与地方实践、报称数字和社会经历分开处理。

\eventanchor{ML-1462-003}{明·1462（天顺六年）}
\paragraph{1462（天顺六年）：景泰、天顺时期的继承、人事与诏令链条可在本年材料中定位，政变责任不可由单一史书裁断。} 这项记录把本章所见的制度运作落实到具体年份。其形成背景是：继承、官僚协调或皇权运作的压力使问题进入朝廷程序。 它带来的、而且仍可由后续材料追踪的改变是：它影响后续人事与政策议程；动机和派系标签不能由单一叙事裁定。 《英宗实录》天顺六年条；《明史·本纪》及相关列传；诏令、奏疏或会典版本 提供相互检验的起点；但桥接条目只标示可回查的年度问题与材料链；实录有中央选择性，崇祯期尤其需要奏疏、地方、朝鲜及满文材料交叉。 因此本段仅据材料确认事件的时序、行动和可见后果，并把政策文本与地方实践、报称数字和社会经历分开处理。

\eventanchor{ML-1462-004}{明·1462（天顺六年）}
\paragraph{1462（天顺六年）：蒙古诸部、朝鲜及西北绿洲的交涉在本年留下多方材料，应以具体使团和地名校正概称。} 这项记录把本章所见的制度运作落实到具体年份。其形成背景是：朝贡名义、边境安全与港口贸易的利益使交涉持续发生。 它带来的、而且仍可由后续材料追踪的改变是：它为后续册封、互市或海防安排留下文书与跨政权比较线索。 《英宗实录》天顺六年条；《明史·外国传》；《朝鲜王朝实录》、琉球《历代宝案》或海外档案 提供相互检验的起点；但桥接条目只标示可回查的年度问题与材料链；实录有中央选择性，崇祯期尤其需要奏疏、地方、朝鲜及满文材料交叉。 因此本段仅据材料确认事件的时序、行动和可见后果，并把政策文本与地方实践、报称数字和社会经历分开处理。

\eventanchor{ML-1463-001}{明·1463（天顺七年）}
\paragraph{1463（天顺七年）：蒙古诸部、朝鲜及西北绿洲的交涉在本年留下多方材料，应以具体使团和地名校正概称。} 这项记录把本章所见的制度运作落实到具体年份。其形成背景是：朝贡名义、边境安全与港口贸易的利益使交涉持续发生。 它带来的、而且仍可由后续材料追踪的改变是：它为后续册封、互市或海防安排留下文书与跨政权比较线索。 《英宗实录》天顺七年条；《明史·外国传》；《朝鲜王朝实录》、琉球《历代宝案》或海外档案 提供相互检验的起点；但桥接条目只标示可回查的年度问题与材料链；实录有中央选择性，崇祯期尤其需要奏疏、地方、朝鲜及满文材料交叉。 因此本段仅据材料确认事件的时序、行动和可见后果，并把政策文本与地方实践、报称数字和社会经历分开处理。

\eventanchor{ML-1463-002}{明·1463（天顺七年）}
\paragraph{1463（天顺七年）：土木危机后的边镇整顿、京师防务或复辟前后军政调度在本年材料中构成连续问题。} 这项记录把本章所见的制度运作落实到具体年份。其形成背景是：前期边防、地方武装或跨区域资源调动的压力推动了该行动。 它带来的、而且仍可由后续材料追踪的改变是：它改变了后续调兵、补给或地方秩序的条件；战果和伤亡须分开核对。 《英宗实录》天顺七年条；《明史·兵志》及相关列传；边镇、地方志或对方材料 提供相互检验的起点；但桥接条目只标示可回查的年度问题与材料链；实录有中央选择性，崇祯期尤其需要奏疏、地方、朝鲜及满文材料交叉。 因此本段仅据材料确认事件的时序、行动和可见后果，并把政策文本与地方实践、报称数字和社会经历分开处理。

\subsection{1463年前后的续接}

\eventanchor{ML-1463-003}{明·1463（天顺七年）}
\paragraph{1463（天顺七年）：军饷、漕运和户籍赋役的本年记录揭示恢复秩序的行政成本，账面额与实收须分列。} 这项记录把本章所见的制度运作落实到具体年份。其形成背景是：财政征解、交通工程或货币支付的压力使相关措施进入政务。 它带来的、而且仍可由后续材料追踪的改变是：其法定安排不等于地方执行，后续须以地域材料追踪负担与成效。 《英宗实录》天顺七年条；《明会典》相关条目；地方志、黄册/契约/河工材料 提供相互检验的起点；但桥接条目只标示可回查的年度问题与材料链；实录有中央选择性，崇祯期尤其需要奏疏、地方、朝鲜及满文材料交叉。 因此本段仅据材料确认事件的时序、行动和可见后果，并把政策文本与地方实践、报称数字和社会经历分开处理。

\eventanchor{ML-1463-004}{明·1463（天顺七年）}
\paragraph{1463（天顺七年）：科举、学校和法律条文在本年制度材料中可追踪，不能据规范文本推定州县实践一致。} 这项记录把本章所见的制度运作落实到具体年份。其形成背景是：制度、教育或知识传播的需要推动了相关行动或文本形成。 它带来的、而且仍可由后续材料追踪的改变是：它提供制度和传播史的锚点，不能据此推断社会整体接受。 《英宗实录》天顺七年条；《明史·选举志》或《艺文志》；文集、碑刻或书目材料 提供相互检验的起点；但桥接条目只标示可回查的年度问题与材料链；实录有中央选择性，崇祯期尤其需要奏疏、地方、朝鲜及满文材料交叉。 因此本段仅据材料确认事件的时序、行动和可见后果，并把政策文本与地方实践、报称数字和社会经历分开处理。

\eventanchor{ML-1464-001}{明·1464（天顺八年）二月23日}
\paragraph{1464（天顺八年）二月23日：天顺帝驾崩，朱千深即位成化帝} 这项记录把本章所见的制度运作落实到具体年份。其形成背景是：继承、官僚协调或皇权运作的压力使问题进入朝廷程序。 它带来的、而且仍可由后续材料追踪的改变是：它影响后续人事与政策议程；动机和派系标签不能由单一叙事裁定。 《英宗实录》天顺八年条；《明史·本纪》及相关列传；诏令、奏疏或会典版本 提供相互检验的起点；但译写候选以编年材料定位；实录记载反映朝廷选择，崇祯期须另以奏疏、地方及敌对政权材料交叉。 因此本段仅据材料确认事件的时序、行动和可见后果，并把政策文本与地方实践、报称数字和社会经历分开处理。

\eventanchor{ML-1464-002}{明·1464（天顺八年）}
\paragraph{1464（天顺八年）：广西瑶族起义军侯大狗} 这项记录把本章所见的制度运作落实到具体年份。其形成背景是：继承、官僚协调或皇权运作的压力使问题进入朝廷程序。 它带来的、而且仍可由后续材料追踪的改变是：它影响后续人事与政策议程；动机和派系标签不能由单一叙事裁定。 《英宗实录》天顺八年条；《明史·本纪》及相关列传；诏令、奏疏或会典版本 提供相互检验的起点；但译写候选以编年材料定位；实录记载反映朝廷选择，崇祯期须另以奏疏、地方及敌对政权材料交叉。 因此本段仅据材料确认事件的时序、行动和可见后果，并把政策文本与地方实践、报称数字和社会经历分开处理。

\eventanchor{ML-1464-003}{明·1464（天顺八年）}
\paragraph{1464（天顺八年）：宝藏航行：宝藏航行的文献被刘大侠从兵部档案中移出并销毁，理由是它们是“离奇之事，远非耳目所证的欺骗性夸大”，“三宝出征西海，浪费了数十万的金钱和粮食，而且死亡的人也可以算进”。虽然他带回了许多珍贵的东西，但这对国家有什么好处呢？} 这项记录把本章所见的制度运作落实到具体年份。其形成背景是：前期边防、地方武装或跨区域资源调动的压力推动了该行动。 它带来的、而且仍可由后续材料追踪的改变是：它改变了后续调兵、补给或地方秩序的条件；战果和伤亡须分开核对。 《英宗实录》天顺八年条；《明史·兵志》及相关列传；边镇、地方志或对方材料 提供相互检验的起点；但译写候选以编年材料定位；实录记载反映朝廷选择，崇祯期须另以奏疏、地方及敌对政权材料交叉。 因此本段仅据材料确认事件的时序、行动和可见后果，并把政策文本与地方实践、报称数字和社会经历分开处理。
